\chapter{Wstęp i cel pracy}

Klasyczne silniki cieplne są jednym z najważnieszych osiągnięc fizyki. Zapoczątkowały one termodynamikę, która jest dziś jednym z najważniejszych działów fizyki. Badanie promieniowania cieplnego dało natomiast początek mechanice kwantowej. W pracy opisano zastosowanie jej do silników cieplnych.

Celem pracy magisterskiej jest zapoznanie się z obecną teorią działania kwantowych silników cieplnych pracujących w różnych cyklach termodynamicznych przy użyciu różnych czynników roboczych.

W pierwszym rozdziale przedstawiona została podstawa termodynamiki fenomenologicznej oraz opisano cykle Carnota, Otto oraz Stirlinga w warunkach równowagowych oraz endoodwracalnych. Warunki równowagowe są jednak niemożliwe do realizacji ponieważ mogą wymagać nieskończenie powolnych przemian lub zbiornkików cieplnych o zmiennej temperaturze. Warunki endoodwracalne są bliższe rzeczywistości, gdyż przepływ ciepła zachodzi przy różnicy temperatur między zbiornikiem a czynnikiem roboczym. Dla obu warunków wyprowadzono sprawności tych cykli termodynamicznych.

W drugim rodziale przedstawiono teorię działania kwantowych cykli termodynamicznych, którą wykorzystano do analizy cyklu Carnota w warunkach równowagowych. Cykle Otto i Stirlinga przeanalizowano dla dwóch czynników roboczych: oscylatora harmonicznego oraz układu o spinie \(\inv{2}\). Pierwszy z cykli rozważono w warunkach równowagowych oraz endoodwracalnych, natomiast drugi jedynie w równowagowych. Analiza obydwu cykli wymagała obliczeń numerycznych, których kod zawarty jest w dodatku.
